\subsection*{E226\quad 短偶极天线的辐射角分布、总功率与辐射电阻}
真空中沿$z$轴的短线性电偶极矩$\boldsymbol p(t)=p_0\cos\omega t\,\hat z$，$\omega>0$，$p_0$为实常数；远场$r\gg\lambda=2\pi c/\omega$，忽略材料损耗和高阶多极矩。
(1) 从迟滞电磁势写出主导$1/r$的辐射电磁场。
(2) 求时间平均的角向Poynting通量与方向图。
(3) 对全立体角积分得到总辐射功率，并与Larmor公式交叉验证。
(4) 对电流幅值$I_0$、电偶极有效长度$\ell$满足$p_0=I_0\ell/\omega$的短天线求辐射电阻。
\par\textbf{逐问解析}
(1) 由时变电偶极的迟滞场远区近似，$\boldsymbol E_{\rm rad}=[\ddot{\boldsymbol p}(t-r/c)\times\hat r]\times\hat r/(4\pi\epsilon_0c^2r)$，对轴向偶极其非零横向分量的幅值为$E_{\theta,0}=p_0\omega^2\sin\theta/(4\pi\epsilon_0c^2r)$。辐射磁场为$H_\phi=E_\theta/Z_0$、$Z_0=\sqrt{\mu_0/\epsilon_0}=\mu_0c$，两场同相且正交，近场$1/r^2,1/r^3$项在此尺度被忽略。
(2) 周期平均Poynting径向分量$\langle S_r\rangle=\langle E_\theta H_\phi\rangle=E_{\theta,0}^2/(2Z_0)=p_0^2\omega^4\sin^2\theta/(32\pi^2\epsilon_0c^3r^2)$。故单位立体角的功率$dP/d\Omega=r^2\langle S_r\rangle$与$\sin^2\theta$成正比，轴向$\theta=0,\pi$没有辐射，而赤道方向最强。
(3) 利用$\int_{\mathbb S^2}\sin^2\theta\,d\Omega=8\pi/3$，得到$\overline P=p_0^2\omega^4/(12\pi\epsilon_0c^3)$。若偶极来自单电荷$q$的小幅运动$p_0=qx_0$，瞬时Larmor功率$q^2a^2/(6\pi\epsilon_0c^3)$，其中$a(t)=-\omega^2x_0\cos\omega t$，对$\cos^2$取平均为$1/2$，恰复现上述$\overline P$。
(4) 因$\dot p(t)$为电流，幅值关系$I_0=\omega p_0/\ell$，设等效天线辐射功率$\overline P=I_0^2R_{\rm rad}/2$，解得$R_{\rm rad}=\omega^2\ell^2/(6\pi\epsilon_0c^3)=Z_0(k\ell)^2/(6\pi)=80\pi^2(\ell/\lambda)^2\,\Omega$。只有$k\ell\ll1$、电流分布能用有效长度$\ell$表示时才适用；真实线天线的输入阻抗还含电抗及导体损耗。
\par\medskip
\subsection*{E227\quad 谐振腔微扰理论的广义本征值频移与介电体插入位置}
无耗散、非色散、各向同性介质的封闭理想导体谐振腔满足$\nabla\times\mu^{-1}\nabla\times\boldsymbol E=\omega^2\epsilon(\boldsymbol r)\boldsymbol E$。假设模为非简并、$\mu$保持不变，在小体积$V_p$内介电常数增加$\delta\epsilon(\boldsymbol r)$且扰动足够弱。
(1) 由Maxwell广义Hermitian本征值问题导出一阶频移公式。
(2) 对一维理想腔驻波$E_0(x)=E_a\sin(\pi x/L)$求局部小介电片位置依赖频移。
(3) 若位置从场节点移向反节点，分析符号、最大幅度与电场能量集中区。
(4) 讨论简并模式、色散材料和强扰动为什么不能直接采用该公式。
\par\textbf{逐问解析}
(1) 写$\mathcal A E=\omega^2\epsilon E$，其中$\mathcal A=\nabla\times\mu^{-1}\nabla\times$在理想导体边界下为自伴算符。对$\epsilon\to\epsilon+\delta\epsilon$和$\omega\to\omega+\delta\omega$作一阶变分，左乘未扰动场$E_0^*$积分，利用自伴性消去$\delta E$项，得$2\omega\delta\omega\int\epsilon|E_0|^2dV=-\omega^2\int\delta\epsilon|E_0|^2dV$。所以$\delta\omega/\omega=-\frac12[\int\delta\epsilon|E_0|^2dV]/[\int\epsilon|E_0|^2dV]$。
(2) 设原介质$\epsilon$均匀，横截面积$A$，薄片厚$\ell\ll L$且位于$x=x_0$，其介电增量为常数$\Delta\epsilon$。分母为$\epsilon A E_a^2\int_0^L\sin^2(\pi x/L)dx=\epsilon A E_a^2L/2$；分子主阶约$\Delta\epsilon A E_a^2\ell\sin^2(\pi x_0/L)$，故$\delta\omega/\omega\simeq-(\Delta\epsilon/\epsilon)(\ell/L)\sin^2(\pi x_0/L)$。
(3) 对$\Delta\epsilon>0$，频移恒为非正：插入高介电材料会降低谐振频率。场节点$x_0=0,L$附近主阶频移近零，反节点$x_0=L/2$最大，值为$-(\Delta\epsilon/\epsilon)(\ell/L)$。该位置依赖可用于腔场测绘或介质传感，且对归一化振幅$E_a$不敏感。
(4) 若模式简并，必须在简并模子空间内对$\delta\epsilon$的矩阵元作对角化；若介电常数随频率变化，正确归一化涉及能量密度的$\partial_\omega(\omega\epsilon)$，不能直接套用非色散分母。若$\Delta\epsilon$或片厚过大，场本身会明显重排，需要求解完整边值问题；小扰动公式只提供首阶渐近。
\par\medskip
